% ===== Джерела до розділу 13. Ключі з префіксом ch13: =====
% Позиції ядра (bib-core.tex) тут НЕ дублюються.

% --- APT: концепція та кейси ---
\bibitem{ch13:mandiant} Mandiant. APT1: Exposing One of China's Cyber Espionage Units. Mandiant Intelligence Center Report, 2013.

\bibitem{ch13:stoll} Stoll C. The Cuckoo's Egg: Tracking a Spy Through the Maze of Computer Espionage. New York: Doubleday, 1989. 326~p.

\bibitem{ch13:virvilis} Virvilis N., Gritzalis D. The Big Four --- What We Did Wrong in Advanced Persistent Threat Detection? \emph{2013 International Conference on Availability, Reliability and Security}. Regensburg: IEEE, 2013. P.~248--254.

% --- Полювання за загрозами ---
\bibitem{ch13:sqrrl} Sqrrl Team. A Framework for Cyber Threat Hunting. Sqrrl Data, 2016.

\bibitem{ch13:hunting-maturity} Lee R. M., Bianco D. The Hunting Maturity Model. SANS Institute / Threat Hunting Project, 2016.

\bibitem{ch13:ocm} Lee R. M., Lee R. The Who, What, Where, When, Why and How of Effective Threat Hunting. SANS Institute Reading Room, 2016.

% --- Living-off-the-land та бокове переміщення ---
\bibitem{ch13:lateral} Smiliotopoulos C., Kambourakis G., Barbatsalou K. On the Detection of Lateral Movement through Supervised Machine Learning and an Open-Source Tool to Create Turnkey Datasets. \emph{International Journal of Information Security}. 2023. Vol.~22. P.~1893--1919.

\bibitem{ch13:gtfobins} GTFOBins Project. Curated List of Unix Binaries That Can Be Used to Bypass Local Security Restrictions. URL: \url{https://gtfobins.github.io/}.

% --- Дані та телеметрія для полювання ---
\bibitem{ch13:oasis-car} MITRE. Cyber Analytics Repository (CAR). The MITRE Corporation. URL: \url{https://car.mitre.org/}.

% --- Українські джерела ---
\bibitem{ch13:hunt-ua} Проактивне полювання за кіберзагрозами (threat hunting) як складова сучасної стратегії кіберзахисту. \emph{Кібербезпека: освіта, наука, техніка}. 2024. URL: \url{https://csecurity.kubg.edu.ua/index.php/journal/article/view/773}.

\bibitem{ch13:apt-crit-ua} Виявлення цілеспрямованих атак на об'єкти критичної інфраструктури України: методи та підходи. \emph{Захист інформації}. 2023. Т.~25, №~3.
